% !TEX root = ../main.tex
% ================== RESULTADOS ==================
% Troque os quadros cinza pelos graficos reais (figuras/*.png) e preencha a tabela.
\secao{Resultados e Discussão}

\titulofig{Ensaio de degrau em malha aberta}
\imagem[0.9\linewidth][12.5cm]{figuras/degrau-malha-aberta.png}

\pendente{Descrever o ensaio: duty aplicado, temperatura inicial e final, e o valor de \(\tau\)
obtido (tempo para atingir 63,2\,\% da variação).}

\titulofig{Resposta em malha fechada ao setpoint}
\imagem[0.9\linewidth][12.5cm]{figuras/malha-fechada.png}

\titulofig{Desempenho do controlador PI}
\begin{center}
\begin{tabular}{lc}
  \toprule
  \textbf{Parâmetro} & \textbf{Valor}\\
  \midrule
  Constante de tempo \(\tau\)        & \pendente{-- s}\\
  Ganhos \(K_p\) / \(K_i\)           & \pendente{-- / --}\\
  Setpoint                           & \pendente{-- °C}\\
  Overshoot                          & \pendente{-- \%}\\
  Tempo de acomodação                & \pendente{-- s}\\
  Erro em regime                     & \pendente{\(\pm\) -- °C}\\
  \bottomrule
\end{tabular}
\end{center}

\pendente{Discutir os resultados: o PI manteve a temperatura no setpoint? Houve overshoot? Como
a compensação de bateria afetou a resposta?}
